\subsection{次数}

设 \(\mathbf{P} = \mathbf{A}[(\mathbf{X}_i)_{i\in I}]\) 是一个多项式代数。对于每个整数 \(n\in \mathbf{N}\) ，令 \(\mathbf{P}_n\) 为 \(\mathbf{P}\) 的由单项式 \(\mathbf{X}^{\nu}\) （其中 \(|\nu| = \sum_{i\in I}\nu_i\)

等于 \(n\) 者）生成的子模。于是 \((\mathbf{P}_n)_{n \in \mathbf{N}}\) 是一个分次，它把 \(\mathbf{A}[(\mathbf{X}_i)_{i \in I}]\) 变成 \(\mathbf{N}\) 型分次代数（III，第459页）。次数为 \(n\) 的齐次元在 \(\mathbf{A}[(\mathbf{X}_i)_{i \in I}]\) 中有时称为 \(n\) 次型（关于未定元 \(\mathbf{X}_i\) ）。

当我们处理非齐次多项式的次数时，我们约定在自然数集 N 中添加一个元素，记作 \(-\infty\) ，并按下列约定把 N 上的序关系和加法扩张到 \(N \cup \{-\infty\}\) （其中 \(n \in N\) ）：

\[
- \infty <   n, (- \infty) + n = n + (- \infty) = - \infty , (- \infty) + (- \infty) = - \infty .
\]

设 \(u = \sum_{\nu \in \mathbf{N}^{(I)}} \alpha_{\nu} X^{\nu}\) 是一个多项式。齐次分量 \(u_{n}\) 的次数为

\(n\) （它是 \(u\) 的齐次分量，对于上述定义的 \(\mathbf{N}\) 型分次而言）等于 \(\sum_{|\nu| = n} \alpha_{\nu} X^{\nu}\) ，并且我们

显然有 \(u = \sum_{n \in \mathbb{N}} u_n\) 。若 \(u \neq 0\) ，这些 \(u_n\) 并非全为零，我们定义 u 的次数

（或全次数），记作 \(\deg u\) ，为使得 \(u_{n} \neq 0\) 的那些 n 中的最大者；换言之（III，第453页），u 的次数是整数 \(|\nu|\) （其中多重指标 \(\nu\) 满足 \(\alpha_{\nu} \neq 0\) ）中的最大者。当 u = 0 时，按约定 u 的次数为 \(-\infty\) 。对每个整数 \(p \in N\) ，关系式 \(\deg u \leqslant p\) 便等价于 \(\alpha_{\nu} = 0\) 对每个多重指标 \(\nu\) （ \(|\nu| > p\) ）成立；因此，满足 \(\deg u \leqslant p\) 的多项式 u 的集合是 \(\mathbf{A}[(\mathbf{X}_{i})_{i \in I}]\) 的一个 A-子模，用上述记号即等于 \(P_{0} + P_{1} + \cdots + P_{p}\) 。

设 E 为一个 A-模。族 \((\mathrm{E} \otimes \mathrm{P}_{n})_{n \in \mathbb{N}}\) 是系数在 E 中的多项式所成的模 \(\mathrm{E}[(\mathrm{X}_{i})_{i \in I}] = \mathrm{E} \otimes_{\mathbb{A}} \mathrm{A}[(\mathrm{X}_{i})_{i \in I}]\) 的一个 N 型分次。我们把上文对非齐次多项式次数所采用的约定推广到这种情形。

命题 1. --- 设 u 和 v 是两个多项式。

\begin{enumerate}
\def\labelenumi{(\roman{enumi})}
\tightlist
\item
  如果 \(\deg u \neq \deg v\) ，则我们有
\end{enumerate}

\[
u + v \neq 0 \quad a n d \quad \deg (u + v) = \sup (\deg u, \deg v).
\]

如果 \(\deg u = \deg v\) ，则我们有 \(\deg (u + v) \leqslant \deg u\) 。

\begin{enumerate}
\def\labelenumi{(\roman{enumi})}
\setcounter{enumi}{1}
\tightlist
\item
  我们有 \(\deg(uv) \leqslant \deg u + \deg v\) .
\end{enumerate}

证明是直接的。

设 \(J \subset I\) ，且 \(B = A[(X_i)_{i \in I - J}]\) ；我们将把 \(A[(X_i)_{i \in I}]\) 与 \(B[(X_i)_{i \in J}]\) 等同（III，第 453 页及以下）。 \(u \in A[(X_i)_{i \in I}]\) 作为 \(B[(X_i)_{i \in J}]\) 的元素的次数，称为 \(u\) 关于 \(X_i\) （指标 \(i \in J\) ）的次数（III，第 454 页）。

设 \(u = \sum_{k=0}^{n} \alpha_k X^k \in A[X]\) 为次数为 \(n\) 的非零多项式，属于单个

不定元。系数 \(\alpha_{n}\) （由假设它 \(\neq0\) ）称为 u 的首项系数。首项系数等于 1 的多项式 \(u\neq0\) 称为首一多项式。在

 \(\mathbf{A}[\mathbf{X}_1, \mathbf{X}_2, \ldots, \mathbf{X}_q]\) 中，总次数为 \(p\) 的单项式的个数，等于形如 \((n_k)_{1 \leqslant k \leqslant q}\) 、属于 \(\mathbf{N}^q\) 且满足 \(\sum_{k=1}^{q} n_k = p\) 的元素的个数，即 \(\binom{q + p - 1}{p}\)

（集合 III，命题 15，第 182 页）。

更一般地，设 \(\Delta\) 是一个交换幺半群， \((\delta_i)_{i \in I}\) 是 \(\Delta\) 的一个元素族。存在唯一的 \(\Delta\) 型分次（作用于代数 \(\mathbf{A}[(\mathbf{X}_i)_{i \in I}]\) ），使得每个单项式 \(\mathbf{X}^\nu\) 的次数为 \(\sum_{i \in I} \nu_i \delta_i\) （III，第458页，

例3）。上述考虑的情形是 \(\Delta = \mathbf{N}\) 且 \(\delta_{i} = 1\) 的情形。在一般情形，为避免混淆，我们将用「权」一词代替「次数」，用「等权」代替「齐次」。例如，存在唯一的 \(\mathbf{N}\) 型分次（作用于代数 \(\mathbf{A}[(\mathbf{X}_{i})_{i\geqslant 1}]\) ），使得 \(\mathbf{X}_i\) 的权为 \(i\) （对每个整数 \(i \geqslant 1\) ）。权为 \(n\) 的等权元素是形如 \(\sum_{\nu} a_{\nu} X^{\nu}\) 的多项式，其中 \(a_{\nu} = 0\) （对于 \(\sum_{i \geqslant 1} i \cdot \nu_i \neq n\) ）。

